\documentclass{article}
\usepackage[english]{babel}
\usepackage[toc,page]{appendix}
\usepackage[letterpaper,top=2cm,bottom=2cm,left=3cm,right=3cm,marginparwidth=1.75cm]{geometry}
\usepackage{enumitem}
\usepackage{amsmath}
\usepackage[colorlinks=true, allcolors=blue]{hyperref}
\usepackage{minted}
\usepackage{float}
\usepackage{graphicx}
\usepackage{multirow}
\usepackage{subcaption}

\title{Title\\[0.3em] \large Report on Assignment 3}
\author{Aaron Bussche \and Jason Kasdiran \and David Magaram}
\date{Group 28}

\begin{document}
\maketitle

\section{3.1}

\section{3.2}
Follow the README to see how to get images and quantitative data. 

Alignment --> Each boid tries to match the vector (speed and direction) of the other boids around it.
Coherence controls how quickly they try to match vectors. 


Coherence --> Each boid tries to move towards the center of mass of the other boids around it. Coherence controls how quickly they try to move towards the center of mass.

Separation --> Each boid tries to keep a certain distance from the other boids around it to not run into them. If boids get too close, they will steer away. Separation controls how much they try to keep this distance.

Visual range --> The distance within which a boid can see other boids and be influenced by them. If the visual range is too small, the boids will not be able to see enough other boids to form a cohesive flock. If it is too large, they may be influenced by too many boids and not form a cohesive flock.

\subsection{Parameter values}



\begin{table}[h!]
    \centering
    \label{tab:params}
    \begin{tabular}{lccc}
        \hline
        \textbf{Parameter} & \textbf{Value}\\
        \hline
        $Radius_{\text{inner}}$ & 10 \\
        $Radius_{\text{outer}}$ & 25 \\
        Alignment weight & 1 \\
        Coherence weight & 1 \\
        Separation weight & 1 \\
        Number of boids & 200 \\
        \hline
    \end{tabular}
    \caption{Parameters.}
\end{table}

\subsection{Results}
\subsection{Visualization}
\subsection{Quantitative data}







\end{document}